% Cover letter and point-by-point reviewer responses.
% Reviewer text is verbatim from plosgen_review.txt; \reply{...} blocks
% are scaffolded as TODO placeholders to be filled in during the revision.
% Style adapted from the mkado-paper review-responses.tex
% (https://github.com/andrewkern/mkado-paper).
%
% This file is a body fragment: it is \input by review-responses-main.tex,
% which supplies the preamble and % Macros for reviewer responses and back-links from response document
% to manuscript line numbers. Pattern adapted from popgensbi manuscript.
%
% Expects the lineno package to be loaded by the calling document and
% \linenumbers to be active in the manuscript body.

% Reviewer / point counters
\newcommand{\thereviewer}{}
\newcounter{point}
\setcounter{point}{0}

% Start a new reviewer block. Argument is the reviewer label
% (e.g. \reviewersection{1} or \reviewersection{AE}).
\newcommand{\reviewersection}[1]{%
  \renewcommand{\thereviewer}{#1}%
  \setcounter{point}{0}%
  \section*{Reviewer \thereviewer:}%
}

% A single review point. Argument is an optional name (often empty).
% Pattern from
% http://tex.stackexchange.com/questions/2317/latex-style-or-macro-for-detailed-response-to-referee-report
\newenvironment{point}[1]
  {\color{gray}\refstepcounter{point}\bigskip\hrule\medskip\noindent
   \slshape{\fontseries{b}\selectfont (\thereviewer.\thepoint) #1}}
  {}

\newcommand{\reply}{\normalfont\medskip\noindent\textbf{Reply}:\ }

% \revpoint{R}{P} marks a manuscript line that addresses reviewer R, point P.
% Use this in the manuscript wherever you make a change in response to a
% reviewer point so the response document can back-link to it.
\newcommand{\revpoint}[2]{\hypertarget{llineno:rev#1:#2}{\linelabel{rr:rev#1:#2}}}

% \revreffull{R}{P} renders a (page, line) reference to such a target.
\newcommand{\revreffull}[2]{{(p.\ \hyperlink{llineno:rev#1:#2}{\pageref{rr:rev#1:#2}, l.\ \lineref{rr:rev#1:#2}})}}

% \revref expands to the page/line reference for the current reviewer/point
% inside a {point} environment.
\newcommand{\revref}{\revreffull{\thereviewer}{\thepoint}\xspace}

% Named back-links: place \llabel{name} in the manuscript, refer to it
% from anywhere with \llname{name}.
\newcommand{\llabel}[1]{\hypertarget{ll:#1}{\linelabel{#1}}}
\newcommand{\llname}[1]{{(p.\ \hyperlink{ll:#1}{\pageref{#1}, l.\ \lineref{#1}})}}

% Italics-quoted block (used to set off the verbatim reviewer text).
\newenvironment{itquote}
  {\begin{quote}\color{gray}\itshape}
  {\end{quote}}

% Convenience macros
\newcommand{\rollover}{\reply{The reviewer makes an excellent point that we have missed out entirely. We have made all the changes suggested, down to the minutiae \revref.}}
\newcommand{\playdead}{\reply{The reviewer makes an excellent point. We have made an utterly trivial change \revref{}that we think deals entirely with the concern raised.}}

% lineno + amsmath fix: lineno doesn't number lines inside align/equation
% environments by default. Patch from
% http://tex.stackexchange.com/questions/43648
\ifcsname patchAmsMathEnvironmentForLineno\endcsname
\else
\newcommand*\patchAmsMathEnvironmentForLineno[1]{%
  \expandafter\let\csname old#1\expandafter\endcsname\csname #1\endcsname
  \expandafter\let\csname oldend#1\expandafter\endcsname\csname end#1\endcsname
  \renewenvironment{#1}%
    {\linenomath\csname old#1\endcsname}%
    {\csname oldend#1\endcsname\endlinenomath}}
\newcommand*\patchBothAmsMathEnvironmentsForLineno[1]{%
  \patchAmsMathEnvironmentForLineno{#1}%
  \patchAmsMathEnvironmentForLineno{#1*}}
\AtBeginDocument{%
  \patchBothAmsMathEnvironmentsForLineno{equation}%
  \patchBothAmsMathEnvironmentsForLineno{align}%
  \patchBothAmsMathEnvironmentsForLineno{flalign}%
  \patchBothAmsMathEnvironmentsForLineno{alignat}%
  \patchBothAmsMathEnvironmentsForLineno{gather}%
  \patchBothAmsMathEnvironmentsForLineno{multline}%
}
\fi
.

\begin{minipage}[b]{2.6in}
  Resubmission Cover Letter \\
  {\it PLOS Genetics} \\
  Manuscript PGENETICS-D-26-00481
\end{minipage}
\hfill
\begin{minipage}[b]{3.4in}
  Angel G.\ Rivera-Col\'on (on behalf of all authors) \\
  \today
\end{minipage}

\vskip 2em

\noindent
{\bf Dear Dr.\ Bergland,}

\vskip 1em

TODO

\vskip 2em

\noindent\hspace{4em}
\begin{minipage}{4.5in}
\noindent
{\bf Sincerely,}

\vskip 2em

{\bf
  Angel G.\ Rivera-Col\'on (on behalf of all authors)
}\\
\end{minipage}

\vskip 4em

\pagebreak

%%%%%%%%%%%%%%%%%%%%%%%%%%%%%%%%%%%%%%%%%%%%%%%%%%%%%%%%%%%%%%%%%%%%%%%%%
\reviewersection{AE}

\begin{itquote}
This manuscript has been reviewed by three expert reviewers as well as
myself. We appreciated the thoroughness of the analysis, and generally
found the writing clear. There was also a shared sentiment that (1) the
focus on genome assembly of hyperpolymorphic species was put forward as
a central aspect of the paper, but there was no assessment presented
that this was the ``best way to do the assembly''; (2) the section on
allozymes fell short because of the limited sampling of individuals (and
no mention of intentional sampling across tidal height suggests that the
relevant haplotypes may have not been sampled); (3) there are multiple
places where the presentation of results could be improved; and (4)
there was insufficient or inaccurate representations of the published
literature. I agree with all of the reviewer's comments. Although the
paper is strong technically, the conceptual advances presented by the
authors are not strong enough for PLOS Genetics. We encourage
resubmission to PLOS Genetics if the authors can incorporate additional
spatial sampling/sequencing across broad (latitudinal) or local
(intertidal height) scales.
\end{itquote}

\reply{%
  \textsc{TODO} --- overall response to the editor.
}

%---------- AE.1 ----------
\begin{point}{} % AE.1 -- Nunez et al. 2021 mischaracterized
Lines 95--98, the authors cite Nunez et al.\ 2021 as an example of a
study focusing on \emph{Mpi}, and follow with a statement that ``this
work has been restricted to a handful of loci\ldots'' That is not a
fair reading of Nunez 2021, which conducted whole-genome pooled
resequencing. To be fair, the reference genome they used was not perfect
but I think that Nunez 2021 does more than you give it credit for.
\end{point}

\reply{%
  \textsc{TODO} --- author response (cite the manuscript change with \texttt{\textbackslash revref}).
}

%---------- AE.2 ----------
\begin{point}{} % AE.2 -- speculative gene function in Results
Lines 228--236. Although the function of these genes may not be
speculative, the discussion of their function as it relates to positive
selection is highly speculative and does not seem relevant for the
Results section.
\end{point}

\reply{%
  \textsc{TODO} --- author response.
}

%---------- AE.3 ----------
\begin{point}{} % AE.3 -- topic sentence at line 271 / trim function speculation
Line 271. This sentence is not a good topic sentence for the paragraph
as it jumps right into the tubular gene and discusses it as the single
example of excess polymorphism. After discussing this gene, the rest of
the paragraph is about genes showing evidence of positive selection. I
suggest that you retool the topic sentence. Also, to reiterate my
comment above --- the function of these genes is not a result of your
work. It is a speculative interpretation about how the likely function
of these genes is the target of flavors of selection. It is probably not
necessary and could be trimmed to make the paper a bit shorter.
\end{point}

\reply{%
  \textsc{TODO} --- author response.
}

%---------- AE.4 ----------
\begin{point}{} % AE.4 -- p-value correction method (Fig 3 legend, throughout)
Figure 3 legend (and throughout). What method did you use to correct
p-values?
\end{point}

\reply{%
  \textsc{TODO} --- author response.
}

%---------- AE.5 ----------
\begin{point}{} % AE.5 -- parenthetical citation style / PLOS numbering
Parenthetical citations: Please correct the style of the parenthetical
citations (and possibly bibliography) to conform to PLOS G's style.
Citations are neither in alphabetical order nor chronological order;
more than one author are listed for pubs with more than two authors;
sometimes there are initials for first \& middle names, sometimes not. I
think PLOS G has numerically ordered citations.
\end{point}

\reply{%
  \textsc{TODO} --- author response.
}

%---------- AE.6 ----------
\begin{point}{} % AE.6 -- intertidal stratification of Mpi sampling
Line 295--351. The classic story from \emph{Semibalanus} is that
\emph{Mpi} is stratified based on position in the intertidal. You all
collected six barnacles but did not, apparently, specifically stratify
on intertidal height. Therefore, you could have missed an important
allelic class (either due to sampling a single position in the
intertidal, or just chance).
\end{point}

\reply{%
  \textsc{TODO} --- author response.
}

%---------- AE.7 ----------
\begin{point}{} % AE.7 -- Figure 4 / Fig S7 tiny text
Figure 4. The text in this figure is tiny. Consider improving for
readability. Same goes for Fig.\ S7.
\end{point}

\reply{%
  \textsc{TODO} --- author response.
}

%---------- AE.8 ----------
\begin{point}{} % AE.8 -- repeats vs diversity claim (lines 376-379)
Lines 376--379. It is a little hard to assess the veracity of this
claim. For example, there is a correlation coefficient of $-0.181$
between $\pi$ and the proportion of repeats (Fig.\ S11). Is that
coefficient statistically different from zero? I bet it is. And, those
conspicuous dips in $\theta_\pi$ and Watterson's $\theta$ at the center
of each chromosome sure do look like centromeres. Do barnacles have
metacentric chromosomes? Also, what are you calling ``repeats'' in 6B?
Could that be anything from a microsat to a TE? How does this
correlation change if you focus only on TEs?
\end{point}

\reply{%
  \textsc{TODO} --- author response.
}

%---------- AE.9 ----------
\begin{point}{} % AE.9 -- "see below" dangling reference
Line 379: ``see below''. See below where?
\end{point}

\reply{%
  \textsc{TODO} --- author response.
}

%---------- AE.10 ----------
\begin{point}{} % AE.10 -- assembly-methods discussion section length
Lines 424--462. I found this section challenging as the first section of
the discussion. You don't really evaluate different methods / approaches
for genome assembly in a highly heterozygous species. Consider reducing
length.
\end{point}

\reply{%
  \textsc{TODO} --- author response.
}

\pagebreak

%%%%%%%%%%%%%%%%%%%%%%%%%%%%%%%%%%%%%%%%%%%%%%%%%%%%%%%%%%%%%%%%%%%%%%%%%
\reviewersection{1}

\begin{itquote}
The authors generated a chromosome-level genome assembly for
\emph{Balanus glandula}, an ecologically important and well-studied
barnacle species in West Coast rocky shore ecosystems. This genome
assembly is of high quality and will be a useful resource for future
evolutionary studies on this species. Additionally, analyses of this
genome assembly, a draft assembly of a congener, and WGS data of 6
individuals from one population, provide insights into the evolutionary
dynamics of species with large population sizes and highly heterozygous
genomes. Overall, the manuscript is well-written, and the methodology is
thorough. However, I have a handful of comments about the manuscript,
which will help ground these analyses in greater taxonomic and
ecological context.
\end{itquote}

\reply{%
  \textsc{TODO} --- overall response to Reviewer~1.
}

%---------- R1.1 ----------
\begin{point}{} % 1.1 -- life history details in the Introduction
In the Introduction there should be increased information about the life
history of the species since these details are pertinent for
understanding demographic and selection processes in the species. For
example, how long is the planktonic larval duration? What is their mode
of reproduction? What is the timing of spawning? How many offspring are
produced per individual?
\end{point}

\reply{%
  \textsc{TODO} --- author response.
}

%---------- R1.2 ----------
\begin{point}{} % 1.2 -- spatial/temporal variation, linear coastline, dispersal
This manuscript focuses extensively on the fact that \emph{B.\ glandula}
is a species with large population sizes. However, there is a large
amount of spatial and temporal variation in \emph{B.\ glandula}
reproductive output, recruitment, and abundance along the west coast,
across both small and large scales (e.g., Gaines, Brown \& Roughgarden
1985; Berger 2009; Broitman et al.\ 2008; MARINe long-term monitoring
data). Additionally, the coastline is linear in nature creating unique
patterns of dispersal and gene flow. These dynamics greatly impact how
evolutionary processes function in this species --- e.g., there is a low
likelihood of adaptation to local conditions since planktonic larvae
from one population often disperses to another population. These dynamics
need to be noted in the Introduction and there should be greater
interpretation of what the genomic patterns of high heterozygosity,
strong positive selection, etc.\ mean within the context of the life
history and ecology of this particular species.
\end{point}

\reply{%
  \textsc{TODO} --- author response.
}

%---------- R1.3 ----------
\begin{point}{} % 1.3 -- novelty of assembly methodology in Results
Given that some of the motivation in the Introduction focuses on the
difficulty in assembling genomes for species with high heterozygosity,
more information should be provided in the results of the genome assembly
to indicate the novelty of the methodology used to produce a very
complete assembly with low duplicates.
\end{point}

\reply{%
  \textsc{TODO} --- author response.
}

%---------- R1.4 ----------
\begin{point}{} % 1.4 -- microhabitat variation and sampling design
More information is needed about the sampling of the individuals from
the Oregon population. It is stated in the Discussion that ``intertidal
microhabitat is approximately uniform in our central Oregon sampling''.
The central Oregon rocky coastline, including Cape Perpetua where the
samples were collected for the WGS, is a dynamic environment with high
levels microhabitat variation (e.g., Sanford \& Menge 2001 provides
evidence that variation in wave exposure within a site influences
post-settlement mortality of barnacle and growth rate). Also, given that
previous research on barnacles has shown that variation in allozyme loci
are associated with spatial heterogeneity (Hedgecock 1994, Schmidt \&
Rand 1999; Nunez et al.\ 2020), why was microhabitat not considered prior
to sampling, as lack of this information greatly weakens the
interpretability of patterns from the data.
\end{point}

\reply{%
  \textsc{TODO} --- author response.
}

%---------- R1.5 ----------
\begin{point}{} % 1.5 -- main-text table of genome statistics
Since a large focus of this manuscript is the generation of the genome
assembly for \emph{B.\ glandula}, the manuscript would benefit from a
table in the main text that summarizes the key genome statistics of the
final assembly --- total length, N50, BUSCO completeness values, etc.\
--- values that are currently in Table S1.
\end{point}

\reply{%
  \textsc{TODO} --- author response.
}

%---------- R1.6 ----------
\begin{point}{} % 1.6 -- B. glandula in middle row of synteny plot (Fig 1)
Since the focus of the manuscript is on \emph{B.\ glandula}, for the
synteny plot in Figure 1, \emph{B.\ glandula} should be in the middle
row, since it would be easier to see the associations of \emph{B.\
glandula} with the other two barnacle species.
\end{point}

\reply{%
  \textsc{TODO} --- author response.
}

%---------- R1.7 ----------
\begin{point}{} % 1.7 -- gene names on Figure 3
To improve the utility of figure 3 in conveying information, it could be
useful to state the gene names on Figure 3 that are mentioned in the
text.
\end{point}

\reply{%
  \textsc{TODO} --- author response.
}

%---------- R1.8 ----------
\begin{point}{} % 1.8 -- Figure 4 axes and legend legibility
The axes and legend for Figure 4 should be bigger to make it easier to
read.
\end{point}

\reply{%
  \textsc{TODO} --- author response.
}

%---------- R1.9 ----------
\begin{point}{} % 1.9 -- congeners terminology and Gpi species attribution
Congeneric barnacles are barnacles belonging to the same genus. Thus,
\emph{Semibalanus balanoides} and \emph{Balanus glandula} are not congers
(line 582). Also, the patterns of variation in \emph{Gpi} identified in
Flight \& Rand 2012 and stated in the following sentence were identified
in \emph{S.\ balanoides}. Thus, it would be more appropriate to state the
species name rather than mention `other marine invertebrates', as that
insinuates that patterns in \emph{Gpi} have been identified more broadly
across marine invertebrate taxa.
\end{point}

\reply{%
  \textsc{TODO} --- author response.
}

%---------- R1.10 ----------
\begin{point}{} % 1.10 -- "small propagule size" (line 612)
It is stated in the Discussion that \emph{B.\ glandula} has ``small
propagule size'' (line 612). \emph{B.\ glandula} is a simultaneous
hermaphrodite that broods its offspring, and releases about 30,000 larvae
per brood (Newman \& Abbott 1980). Thus, similar to many other marine
invertebrates that broadcast spawn, the species has high fecundity and
high levels of early life mortality.
\end{point}

\reply{%
  \textsc{TODO} --- author response.
}

%---------- R1.11 ----------
\begin{point}{} % 1.11 -- methods thorough (positive)
The methods are incredibly thorough and easy to follow.
\end{point}

\reply{%
  \textsc{TODO} --- author response.
}

%---------- R1.12 ----------
\begin{point}{} % 1.12 -- broken code link, data preview for reviewers
The link for the code is not working. Also, the final assembly and
annotation files should be available for preview to reviewers.
\end{point}

\reply{%
  \textsc{TODO} --- author response.
}

%---------- R1.13 ----------
\begin{point}{} % 1.13 -- Line 205: repeated word "consistent"
Line 205: For clarity in reading, change one of the words ``consistent''
to something else, since it is used twice within the same sentence.
\end{point}

\reply{%
  \textsc{TODO} --- author response.
}

%---------- R1.14 ----------
\begin{point}{} % 1.14 -- Line 234: citation initials error (Cheng et al. 2021)
Line 234: Error in the initials for the citation C.H.\ Cheng et al.\
2021.
\end{point}

\reply{%
  \textsc{TODO} --- author response.
}

\pagebreak

%%%%%%%%%%%%%%%%%%%%%%%%%%%%%%%%%%%%%%%%%%%%%%%%%%%%%%%%%%%%%%%%%%%%%%%%%
\reviewersection{2}

\begin{itquote}
This is a very comprehensive study on the genome of a marine invertebrate
with extremely large population size. I commend the authors on putting
together a manuscript that was clear, technically accurate, and has more
depth than an average genome paper. Overall I just have a few
recommendations for improvements.
\end{itquote}

\reply{%
  \textsc{TODO} --- overall response to Reviewer~2.
}

%---------- R2.1 ----------
\begin{point}{} % 2.1 -- length and accessibility
While the study is comprehensive and technically accurate, the manuscript
is very long and gets technically difficult to follow in some places. I
make a few suggestions below, but I encourage the authors to reduce the
length of the manuscript, revise some text to make it more broadly
accessible to users, and tweaking some of the graphics to make them more
intuitive for readers.
\end{point}

\reply{%
  \textsc{TODO} --- author response.
}

%---------- R2.2 ----------
\begin{point}{} % 2.2 -- 6 vs 8 samples ambiguity
It is unclear what the sample size for the population genomics is. Some
places mention 6 samples and some mention 8 samples.
\end{point}

\reply{%
  \textsc{TODO} --- author response.
}

%---------- R2.3 ----------
\begin{point}{} % 2.3 -- cut/shorten allozyme comparison
One section that could be cut from the manuscript is the comparison to
previous allozyme studies because more samples are needed to do this
comparison well. The previous allozyme work was about spatial
heterogeneity and local adaptation within a species, while the tests
being done for selection are more about fixation of different nucleotides
between species. While the latter is a type of test of selection, it
doesn't really get at whether the microgeographic signals reflect
sweepstakes or local adaptation. This analysis is better saved for a more
comprehensive analysis with more samples, or moved to the supplement and
substantially shortened in the main text.
\end{point}

\reply{%
  \textsc{TODO} --- author response.
}

%---------- R2.4 ----------
\begin{point}{} % 2.4 -- present "extremity" earlier with comparative numbers
One of the main themes in the paper is how ``extreme'' this species is in
comparison to other species, but the extremity is not presented well
until the reader gets to the discussion (and even there it is often
presented in relative terms). I recommend adding some levels of
heterozygosity and nucleotide diversity observed in various wild
populations of well-known species in the introduction, and adding some
statistics of other species on the figures and in the Results text to
help readers understand the context.
\end{point}

\reply{%
  \textsc{TODO} --- author response.
}

%---------- R2.5 ----------
\begin{point}{} % 2.5 -- coverage distribution to show haplotigs purged
The statistics presented did not exactly convince me that the genome was
purged of haplotigs. Could you plot a distribution of the coverage to
show that it is not bimodal? (e.g.\
\url{https://doi.org/10.1111/1755-0998.13801})
\end{point}

\reply{%
  \textsc{TODO} --- author response.
}

%---------- R2.6 ----------
\begin{point}{} % 2.6 -- mutation load paradox in large populations
A main take-home is that purifying selection is very efficient (as
expected), but this kind of glosses over the paradox that there are more
deleterious mutations segregating in larger populations and this can lead
to higher load (\url{https://academic.oup.com/cz/article/62/6/567/2745835}).
It would be helpful to but the results of this study in context of that
review in the discussion.
\end{point}

\reply{%
  \textsc{TODO} --- author response.
}

%---------- R2.7 ----------
\begin{point}{} % 2.7 -- explain test statistics (NI, codon bias) in Results
There are a number of statistics that readers may or may not be familiar
with, so it would be helpful to add a short statement about each test
statistic and its biological interpretation in the results. For example
the Neutrality Index is not explained; Figure 3 gets a little confusing
with lower values of NI on the right side. Adding some biological context
(e.g.\ show areas of the plot inferred to be neutral) would help the
results reach a wider audience. As another example, biased use of codon
results could illustrate to the reader the values to expect if there was
no bias. Etc.
\end{point}

\reply{%
  \textsc{TODO} --- author response.
}

%---------- R2.8 ----------
\begin{point}{} % 2.8 -- Ne estimation issues and confidence intervals
Estimating $N_e$ in large populations is statistically problematic with
many estimators (see
\url{https://onlinelibrary.wiley.com/doi/abs/10.1111/faf.12338} and
related papers) and no confidence intervals are presented; it would be
helpful to mention the statistical issues in the discussion.
\end{point}

\reply{%
  \textsc{TODO} --- author response.
}

%---------- R2.9 ----------
\begin{point}{} % 2.9 -- cite reasons synonymous sites under selection
In the discussion it would strengthen the case of Synonymous
substitutions being under selection to cite some reasons why; for example
they can affect the stability of RNAs.
\end{point}

\reply{%
  \textsc{TODO} --- author response.
}

%---------- R2.10 ----------
\begin{point}{} % 2.10 -- dn/ds and MK paradox: cite residue-level examples
In the discussion on the paradox of the dn/ds and McDonald-Kreitman test,
it might be helpful to cite examples of the majority of a coding sequence
can be under purifying selection while specific residues are under
positive selection (e.g.\ \url{https://doi.org/10.1038/nrg3905}).
\end{point}

\reply{%
  \textsc{TODO} --- author response.
}

%---------- R2.11 ----------
\begin{point}{} % 2.11 -- closing remarks (positive)
Overall there were many sections of the paper I enjoyed, especially the
challenges of genome assembly and the paradox of the dn/ds and
McDonald-Kreitman test. It's a fascinating paper and with some revisions
it will make a nice addition to the literature.
\end{point}

\reply{%
  \textsc{TODO} --- author response.
}

\pagebreak

%%%%%%%%%%%%%%%%%%%%%%%%%%%%%%%%%%%%%%%%%%%%%%%%%%%%%%%%%%%%%%%%%%%%%%%%%
\reviewersection{3}

\begin{itquote}
The authors present a new genome assembly for the ecologically important
barnacle \emph{Balanus glandula} and generate genomic data for one
population sample in Oregon and additional assemblies of relatives to
facilitate analyses of variation and selection. I have no major concerns
about the methods used in the paper and appreciate the significant effort
invested to assemble the genome and execute the analyses. The paper is
clear and the methods and supplement are thorough. My main concern is
around the framing of the work. As written, the primary argument seems to
be that we need more studies of organisms with big population sizes and
here we have executed another one --- I think this argument is not the
most compelling one they could make. They could talk more about the
organism/ecology and emphasize what specifically this work adds to the
field (i.e., the exceptional ecological context of this barnacle and its
staggering census size, testing the predictions of extremely effective
purifying and positive selection). Throughout the paper, the authors also
bring up interesting challenges in this area (Lewontin's Paradox, dealing
with hyperpolymorphic genomic data) but don't really squarely address
them. The outcome is that the major takeaways (which are interesting) end
up falling a bit flat. More comprehensively comparing and contrasting
results in this barnacle with other studies of large populations would
also help highlight what this work adds (even better with a figure).
Potentially simulations exploring the implications of the results would
add to the work as well (i.e., can this help explain strong clines over
small spatial scales despite high dispersal? How much more effective
should selection be in this system vs flies? What does these results
predict about response to changing conditions?). A more minor concern ---
I think organizing the paper such that the section on diversity in the
assembly is followed by the sections on diversity in populations and on
estimating effective population size (before getting into HCE's and
selection) would make for a clearer conceptual thread. Predictions about
selection make more sense after getting the info on variation and
effective population size and it is awkward skipping ahead in multiple
places (e.g., 240--241) where the manuscript has to allude to a section
that hasn't come yet.
\end{itquote}

\reply{%
  \textsc{TODO} --- overall response to Reviewer~3.
}

%---------- R3.1 ----------
\begin{point}{} % 3.1 -- Line 53: "few key exceptions" framing
Line 53. Unless the authors are certain they have comprehensively
assessed the literature for empirical studies of variation and selection
in large population sizes, it would be better to phrase this sentence
differently along the lines of ``while there have been empirical studies
on these mechanisms in large populations (e.g.\ \ldots\ldots), they
remain rare.'' The ``few key exceptions'' is problematic if some have been
left out --- for example, I can think of Halligan et al.\ 2010 PLOS
Genetics which looks at wild house mice with a large effective population
size (and census size) as well as Drosophila. I think it is fine to argue
that this topic is understudied as a motivation for the work, but it is
not the strongest argument. Adding one taxon to a big gap is not
necessarily that compelling. However, emphasizing the unique aspects of
this system (the truly enormous census sizes, interesting examples of the
ecological context, clines despite significant potential for dispersal,
etc.)\ and what we can gain from this specific work would be more
persuasive.
\end{point}

\reply{%
  \textsc{TODO} --- author response.
}

%---------- R3.2 ----------
\begin{point}{} % 3.2 -- Line 96: Alm Rosenblad et al. 2021 uses genomic data
Line 96- this seems incorrect given the citation of Alm Rosenblad et al.\
2021 which does use genomic data to estimate diversity.
\end{point}

\reply{%
  \textsc{TODO} --- author response.
}

%---------- R3.3 ----------
\begin{point}{} % 3.3 -- Line 152: phylogeny to interpret result
Line 152. It would be helpful to have a phylogeny to help interpret this
result even if it is a cartoon type tree or in the supplement.
\end{point}

\reply{%
  \textsc{TODO} --- author response.
}

%---------- R3.4 ----------
\begin{point}{} % 3.4 -- Line 165: heterozygosity table by site class
Line 165- A table of expected/observed heterozygosity per site
classification type (coding, non-coding, 0, 4-fold, etc.)\ is needed.
\end{point}

\reply{%
  \textsc{TODO} --- author response.
}

%---------- R3.5 ----------
\begin{point}{} % 3.5 -- Line 183: define HCE criteria in main text
Line 183- Please define criteria for HCE in the main text.
\end{point}

\reply{%
  \textsc{TODO} --- author response.
}

%---------- R3.6 ----------
\begin{point}{} % 3.6 -- Lines 186-190 unclear
Lines 186--190. This is unclear.
\end{point}

\reply{%
  \textsc{TODO} --- author response.
}

%---------- R3.7 ----------
\begin{point}{} % 3.7 -- Lines 226-228: why unknown function?
Lines 226--228 --- Not sure about this explanation. Is the function
unknown bc of poor annotation or lack of functional annotation in this
species? Or do you mean that the steps you took to curate the set
necessarily restricted the analysis to genes with less info about
function for some reason?
\end{point}

\reply{%
  \textsc{TODO} --- author response.
}

%---------- R3.8 ----------
\begin{point}{} % 3.8 -- Line 239: justify sample of six
Line 239- Justify a sample of six individuals (or explicitly discuss
limitations for both variation and selection).
\end{point}

\reply{%
  \textsc{TODO} --- author response.
}

%---------- R3.9 ----------
\begin{point}{} % 3.9 -- SFS figure
A SFS figure would be helpful.
\end{point}

\reply{%
  \textsc{TODO} --- author response.
}

%---------- R3.10 ----------
\begin{point}{} % 3.10 -- Line 252: reconcile MK/dn-ds tests in place
Line 252 --- why not reconcile this right here? For readers that do know
these tests well, it is unsatisfying to delay because the explanation is
clear. Readers that don't understand these tests will appreciate being
able to interpret the results right here. Alternatively, it could be
effective to include the explanation up front in the framing --- explain
why the tests differ and how they assess different things before you even
give the results. Then you can wrap up this section with a very clear
statement that there is evidence of highly effective positive and
purifying selection.
\end{point}

\reply{%
  \textsc{TODO} --- author response.
}

%---------- R3.11 ----------
\begin{point}{} % 3.11 -- Line 361: long or short read?
Line 361- Is this long or short read?
\end{point}

\reply{%
  \textsc{TODO} --- author response.
}

%---------- R3.12 ----------
\begin{point}{} % 3.12 -- Line 409: theta over which sites; census size placement
Line 409 -was this an average theta over all sites or over putatively
neutral sites (intergenic, 4-fold, etc.)? Also, is there a bound on the
estimate of mutation rate that we should consider in estimating the
effective population size? What exactly is the census size? I saw a
density estimate earlier but not an actual census size estimate. A census
size estimate is given in the next to last part of the paper but put it
here or earlier.
\end{point}

\reply{%
  \textsc{TODO} --- author response.
}

%---------- R3.13 ----------
\begin{point}{} % 3.13 -- Line 422: Lewontin's Paradox unexplained
Line 422-Lewontin's Paradox is mentioned without explanation for the
first time here. It really should be discussed in the framing of the work
or explained more here or not mentioned all until the end of the
discussion in the context of future directions. The mention here is
confusing as it is.
\end{point}

\reply{%
  \textsc{TODO} --- author response.
}

%---------- R3.14 ----------
\begin{point}{} % 3.14 -- Section at 424: goal of assembly-methods section unclear
Section starting at 424. The goal of this section is not clear. The key
point seems to be that genome assembly is challenging for highly
heterozygous genome but that there are many new tools that help and it
would be great to assess them w/r to their potential to address those
challenges. However, this paper doesn't do any of those things so it
seems unsatisfying to bring it up and then just say any new genome
produced is progress?
\end{point}

\reply{%
  \textsc{TODO} --- author response.
}

%---------- R3.15 ----------
\begin{point}{} % 3.15 -- Line 478: reword "functional classes of sites"
Line 478 --- unclear, try ``that the functional classes of sites targeted
by purifying selection\ldots..'' As it reads, it sounds more specific
(like the same genes targeted).
\end{point}

\reply{%
  \textsc{TODO} --- author response.
}

%---------- R3.16 ----------
\begin{point}{} % 3.16 -- Line 498: divide sites and look at constraint by class
Line 498- could you divide these sites up and look at signals of
constraint as a class?
\end{point}

\reply{%
  \textsc{TODO} --- author response.
}

%---------- R3.17 ----------
\begin{point}{} % 3.17 -- Line 500: silencing mechanisms for high TE content
Line 500- Is there evidence of silencing mechanisms as with plants that
also have high TE/repeat content or could that be an explanation?
\end{point}

\reply{%
  \textsc{TODO} --- author response.
}

%---------- R3.18 ----------
\begin{point}{} % 3.18 -- Line 554: more cross-taxa comparisons / figure
Line 554- It would strengthen the paper to have more comparisons like
these with other taxa to put the results in context. There is a great
opportunity for a figure contrasting census size, effective population
size, $\pi$, selection, etc.\ among large population species.
\end{point}

\reply{%
  \textsc{TODO} --- author response.
}

%---------- R3.19 ----------
\begin{point}{} % 3.19 -- Section at 556: develop allozyme biology first
Section at 556- It would be good to better develop the biology and
significance of these allozyme loci before discussing the results. For
example, the sentence at 583 should begin the paragraph. Many people will
not be familiar with the history of work in this field and those that are
will appreciate revisiting them to reflect on whether this work sheds new
light.
\end{point}

\reply{%
  \textsc{TODO} --- author response.
}

%---------- R3.20 ----------
\begin{point}{} % 3.20 -- Line 610: "Lancelets"
Line 610-Lancelets.
\end{point}

\reply{%
  \textsc{TODO} --- author response.
}

%---------- R3.21 ----------
\begin{point}{} % 3.21 -- Line 616: name other hyperpolymorphic systems in intro
Line 616-What are the other systems that are hyperpolymorphic? Maybe
discuss these earlier in the intro.
\end{point}

\reply{%
  \textsc{TODO} --- author response.
}

%---------- R3.22 ----------
\begin{point}{} % 3.22 -- Line 620: "of/above"
Line 620- ``of/above''.
\end{point}

\reply{%
  \textsc{TODO} --- author response.
}

%---------- R3.23 ----------
\begin{point}{} % 3.23 -- Lines 637: "this comes so late!"
Lines 637- this comes so late!
\end{point}

\reply{%
  \textsc{TODO} --- author response.
}

%---------- R3.24 ----------
\begin{point}{} % 3.24 -- Lines 667: sperm or diploid DNA?
Lines 667- I am a little confused about what DNA was extracted. Was it
from sperm or diploid cells?
\end{point}

\reply{%
  \textsc{TODO} --- author response.
}
